\section{Approach}

We illustrate our approach using a running example from the \texttt{kinmail\_srs} document, where we generated the \texttt{user\_registration} functionality.

\subsection{Step 1: SRS Processing and Intermediate Representation Generation}

A user uploads an SRS document in PDF, DOCX, or TXT format. For PDF documents, we extract text content while preserving document structure and formatting. The extracted text is segmented into semantically coherent chunks, where each chunk represents a logical section such as a feature description, use case, or functional requirement. For the \texttt{kinmail\_srs} document, the parser extracted chunks including sections describing user authentication, product management, and user registration workflows.

Each text chunk is processed to identify and extract individual functionalities. For each functionality, the system extracts structured information including the functionality name (a unique identifier such as ``user\_registration''), the context (surrounding requirements and related features), the use case (the specific scenario the functionality addresses), detailed requirements (both functional and non-functional), input/output specifications (expected inputs, outputs, and data flow), and business rules (domain-specific rules and constraints). For the \texttt{user\_registration} functionality, the system extracted the name as \texttt{user\_registration}, identified the context as part of the user authentication system, determined the use case as new users creating accounts with email, username, and password, extracted requirements including email validation and password strength requirements, and identified business rules requiring email uniqueness and password strength criteria. To enable semantic search, we create embeddings for each extracted functionality, which are stored in memory and allow users to search for functionalities using natural language queries.

The extracted functionality information is used to generate an Intermediate Representation (IR), which serves as a structured specification that bridges the gap between natural language requirements and executable code. The IR is a JSON specification that defines the exact structure, naming conventions, and architectural patterns for implementing the functionality. The IR includes functionality metadata (name, description, target programming language, and framework), complete file structure specification (all files to be generated including models, controllers, views, and templates with exact file paths), exact naming conventions (names for functions, classes, variables, and database tables that must be used consistently), database schema definitions (model definitions with field names, types, constraints, and relationships), API structure (route definitions, HTTP methods, function signatures, and parameter specifications), exact import paths (import statements that must be used in both code and tests), and framework configuration (framework-specific setup requirements such as Flask blueprints, Express routes, or Spring Boot annotations). The IR serves as a single source of truth that ensures perfect alignment between generated code and tests. By defining exact naming conventions and file structures before code generation, the IR eliminates the common problem where tests and code diverge due to inconsistent naming or structure. Additionally, the IR enforces architectural consistency across all generated code, ensuring that shared models (such as User, Product, Category) are properly referenced to prevent database table collisions. For the \texttt{user\_registration} example, the generated IR specified models in \texttt{models/user.py} (shared User model) with exact fields including \texttt{id} as Integer primary key, \texttt{email} as String(120) unique and not null, \texttt{username} as String(80) unique and not null, and \texttt{password\_hash} as String(255). The controller was specified as \texttt{controllers/user\_registration\_controller.py} with route \texttt{/register} handling GET and POST methods, with the exact function name \texttt{register}. The IR also specified views in \texttt{views/user\_registration\_views.py}, templates in \texttt{templates/user\_registration\_register.html}, import paths as \texttt{from models.user import User} and \texttt{from controllers.user\_registration\_controller import register}, and Flask blueprint registration in \texttt{app.py}. The IR generation analyzes the extracted functionality information and creates this structured specification, enforcing critical rules where shared models are referenced consistently across functionalities, non-shared models are named based on functionality for uniqueness, and all file paths and import statements must match exactly.

\subsection{Step 2: Test-Driven Code Generation}

Following Test-Driven Development (TDD) principles, we generate test cases \textit{before} generating the implementation code. The IR and the extracted functionality information (including requirements, use cases, and business rules) are used for test generation. This ensures that code is written specifically to satisfy test requirements, leading to better test coverage and more reliable implementations. The test generator creates comprehensive test suites that employ a hybrid validation approach. Functional validation ensures tests validate SRS requirements for functional correctness. Structural validation ensures tests define expected code structure based on the IR, ensuring that generated code matches the exact naming conventions, file paths, and import statements specified in the IR. Edge case coverage ensures tests specify expected behavior for edge cases, error conditions, and boundary scenarios. Integration validation ensures tests verify that components work together correctly, including database operations, API endpoints, and view rendering. The tests are designed to import directly from the code files that will be generated in the next phase, ensuring that the generated code must conform to the IR specifications. This TDD approach guarantees that all generated code has corresponding tests, code is written to pass tests rather than tests being written to match code, and test coverage is comprehensive from the start. For the \texttt{user\_registration} example, the system generated comprehensive test cases covering functional tests including email validation for valid and invalid formats, password strength requirements, duplicate email prevention, and successful registration flow. Structural tests verified correct import from \texttt{controllers.user\_registration\_controller}, exact function name \texttt{register}, and proper route handling. Edge case tests covered empty form submission, password mismatch scenarios, and special characters in username. Integration tests verified database user creation, session management, redirect after registration, and error message display. All tests include proper mocking for external dependencies such as email services, file uploads, and database operations, and explicitly import from \texttt{controllers.user\_registration\_controller} and call the \texttt{register} function, enforcing the IR's naming requirements.

Once tests are generated, the IR and functionality information are used for code generation. The code generator produces complete MVC code that satisfies the test requirements and adheres to the IR specifications exactly. The MVC (Model-View-Controller) architecture provides several benefits including separation of concerns where models handle data and business logic, controllers manage request handling and coordination, and views handle presentation, enabling clear separation of responsibilities. The modular structure improves maintainability by making code easier to understand, modify, and extend. The architecture enhances testability as each component can be tested independently, facilitating comprehensive test coverage. The architecture supports scalability by allowing components to be modified or extended without affecting others, and promotes reusability as models and controllers can be reused across different views and functionalities. The code generator produces data models with proper database relationships, validation logic, and field definitions matching the IR, controllers with business logic and request handling using exact function names and signatures specified in the IR, view logic for processing and formatting data, HTML templates for web applications with proper form handling and error display, and framework setup including initialization code for the chosen framework (Flask, Express, Spring Boot, etc.) with blueprint registration, database configuration, and middleware setup. Since tests are generated first, the code generator receives the test specifications as part of its input, ensuring that the generated code will pass the tests. The code generator also receives the IR, ensuring that all naming conventions, file paths, and structural requirements are followed exactly. For the \texttt{user\_registration} example, the system generated a User model in \texttt{models/user.py} with Flask-SQLAlchemy, including \texttt{email}, \texttt{username}, and \texttt{password\_hash} fields with proper validation, matching the IR specification exactly. The controller in \texttt{controllers/user\_registration\_controller.py} contains the \texttt{register()} function handling GET requests for displaying the form and POST requests for processing registration, with email validation, password hashing, and duplicate email checking, using the exact function name specified in the IR. The template in \texttt{templates/user\_registration\_register.html} provides a registration form with email, username, password, and confirm password fields, with both client-side and server-side validation. The framework setup includes blueprint registration in \texttt{app.py}, database initialization, and Flask-Mail configuration for email verification. The generated code structure exactly matches the IR specification, ensuring that tests can import and execute correctly, with all file paths, function names, and import statements aligning with the IR requirements.

\subsection{Step 3: Test Execution and Automated Feedback Loop}

The generated code is exported to the project directory with proper MVC file structure. The system then executes the generated test cases to validate the code. For Python applications, we use pytest to run the test suite. The system captures comprehensive test results including the total number of tests, number of passing tests, number of failing tests, and detailed error messages for each failure. For the \texttt{user\_registration} example, the initial test execution revealed several issues including missing email validation logic in the registration handler, incorrect password hashing implementation, missing error handling for duplicate email scenarios, template rendering issues with form validation messages, and missing database transaction handling. If all tests pass, the process completes successfully. However, when tests fail, the automated feedback loop begins to correct the code iteratively.

The automated feedback loop performs up to three iterations of code correction based on test failures. For each iteration, the system reads the current code state from disk, ensuring it works with the latest version of all files. All files in the MVC structure (models, controllers, views, templates) are reconstructed into a single code string representation. Test failures and detailed error messages are extracted and analyzed to identify specific issues, with the system categorizing failures by type (e.g., assertion errors, import errors, runtime exceptions) and identifying the affected components. For correction, the system is called with the current code including all files in the MVC structure, test failures with detailed error messages and stack traces, the original IR for structure reference and naming requirements, extracted functionality information for requirements reference, and an enhanced prompt matching the structure of initial code generation. For code validation, the response is validated to ensure file markers are present with at least 3 files generated, code length is at least 70\% of input length to prevent incomplete responses, and MD5 hash comparison detects if identical code is returned to prevent infinite loops. Once validated, the corrected code is exported to the project directory, overwriting previous versions. Tests are then run again on the corrected code. If all tests pass, the loop terminates. If tests still fail and iterations remain, the process continues to the next iteration.

For the \texttt{user\_registration} example, the first iteration read all code files, extracted error messages from the failing tests, and generated corrected code that addressed the duplicate email validation issue by adding proper database query checks. The system validated the response, exported the corrected code, and re-ran tests. The second iteration analyzed the remaining failures, focusing on password hashing and template rendering issues. The system generated code with corrected password hashing and improved template error handling. The final iteration addressed remaining edge cases in form validation and database transaction handling. The system generated code with proper try-except blocks for database operations and improved form validation logic. After three iterations, the feedback loop terminated as it reached the maximum iteration limit. The automated feedback loop demonstrates the system's ability to iteratively improve code quality without human intervention, addressing critical issues including email validation, password hashing, error handling, and database operations.
